
% \section{Examples of RDMA Flagged Entities in Dataset Refinement} \label{appendix: dataset refinement}
% We investigate which entities were primarily flagged by RDMA for human review. We visualize the top cases in Fig. \ref{fig:Top_Corrections}, where we observe that many entities mined by \cite{dong2023ontology_poor_annotations} were indeed not rare diseases, while several key rare disease entities were missed in their original annotations.

% \begin{figure}[h!]
%     \centering
%     \includegraphics[width=1.0\textwidth]{figs/RDMA_Corrections.drawio.png}
%     \caption{\textbf{Top 5 Corrected Annotations with RDMA.} RDMA effectively assists human annotation by identifying errors in the noisy historical dataset from \cite{dong2023ontology_poor_annotations}. The figure highlights the top 5 corrections, including both false negatives (a) such as seizure disorder incorrectly labeled as a rare disease in Orphanet, and false positives (b) that were missed in the initial annotations. Of the 72 annotations flagged by RDMA for human review, 55 (76.4\%) were correctly identified as erroneous, demonstrating its value as an annotation assistant for prioritizing human supervision. While hemochromatosis is typically considered a rare condition, in ICU contexts such as MIMIC notes, our physician deemed that to not necessarily be the case. Furthermore, Orphanet \cite{weinreich2008orphanet} makes a distinction between \textbf{rare} hemochromatosis and hemochromatosis. }
%     \label{fig:Top_Corrections}
% \end{figure}
